\pdfbookmark[1]{Abstract}{Abstract}
\thispagestyle{empty}
\begin{center}
  \rule{\textwidth}{0.4pt}\\[0.9em]
  {\Large\spacedallcaps{Abstract}}\\[0.7em]
  \rule{\textwidth}{0.4pt}
\end{center}
\bigskip

\noindent The orior transpiler moves an operation from one language to another through one intermediate
form. A symbol of a source language maps to a bridge word of the transpiler. The bridge word maps to an anchor
of the ladder, and the anchor is a relation on machine words. The relation's cases are the operation's gsm, the
relational assembly of gnascor. A second table maps the bridge word back out to a symbol of the target language.

This research paper states that transform as a composition of finite maps and proves each of its parts in Lean
4. The proofs cover three faces. The ladder relations hold on every case that ladder.h states, at every width and
in both readings of the sign. The matrices of gsm obey five laws over all 86 of their rows. Each operator of
\texttt{.g}, the high order face, compiles to the gsm of one anchor, for every pair of words of every width.

Each target language joins the bridge through a table and a Lean file of its own. The table names the bridge
word of each symbol. The Lean file proves that the symbol's semantics, by its own standard, is the anchor's
relation. A language check fails where a witnessed symbol lacks its theorem. A lift from one language to another
then preserves meaning by transitivity, and it refuses every symbol that has no proof.
